\section{Preliminaries --- Graph Neural Networks}
\label{sec_meth_gnn}

A human skeleton is a graph:
seventeen keypoints form the nodes and the skeletal connections form the edges.
A graph neural network (GNN) embeds each node while conditioning on the structure of the graph
--- the natural choice for producing pose-grouping embeddings that are aware of joint-type and connectivity information.
This section introduces the message-passing framework that unifies modern GNN variants,
derives the graph convolutional network (GCN) as the canonical spectral instance,
and develops the graph attention network (GAT),
culminating in the GATv2 correction that is used as the attention primitive in every experiment in this chapter.

\subsection{The Message-Passing Framework}
\label{sec_meth_mp}

Let $\mathcal{G} = (\mathcal{V}, \mathcal{E})$ denote an undirected graph
with $N$ nodes and edge set $\mathcal{E} \subseteq \mathcal{V} \times \mathcal{V}$.
Each node $i \in \mathcal{V}$ is associated with a feature vector $\mathbf{h}_i^{(0)} \in \mathbb{R}^{d_0}$,
and each edge $(i,j) \in \mathcal{E}$ optionally carries an edge feature $\mathbf{e}_{ij}$.
The neighbourhood of node $i$ is $\mathcal{N}(i) = \{ j : (i,j) \in \mathcal{E} \}$.

Almost every modern GNN variant can be expressed as iterated message passing~\cite{gilmer2017mpnn}.
At layer $l+1$,
each node receives a message from each of its neighbours,
aggregates them with a permutation-invariant operator,
and updates its representation accordingly.
Formally,
\begin{equation}
    \mathbf{m}_{i \leftarrow j}^{(l+1)} = \mathrm{MSG}^{(l)}\!\left(\mathbf{h}_i^{(l)},\, \mathbf{h}_j^{(l)},\, \mathbf{e}_{ij}\right),
    \label{eq:mpnn_message}
\end{equation}
\begin{equation}
    \mathbf{h}_i^{(l+1)} = \mathrm{UPDATE}^{(l)}\!\left(\mathbf{h}_i^{(l)},\, \bigoplus_{j \in \mathcal{N}(i)} \mathbf{m}_{i \leftarrow j}^{(l+1)}\right),
    \label{eq:mpnn_update}
\end{equation}
where $\bigoplus$ denotes a permutation-invariant aggregator (sum, mean, or max)
and $\mathrm{MSG}^{(l)}$ and $\mathrm{UPDATE}^{(l)}$ are learnable functions at layer $l$.
Different choices of these three components recover different GNN architectures.
GCN (Section~\ref{sec_meth_gcn}) uses a linear message,
a degree-based weighted sum aggregator,
and a linear--nonlinear update.
GAT (Section~\ref{sec_meth_gat}) introduces learned attention weights on the aggregator,
making each neighbour's contribution data-dependent.
GraphSAGE~\cite{hamilton2017graphsage} uses concatenated mean--max aggregation with a feedforward update,
optimised for inductive settings.

The message-passing view clarifies two points that matter for the remainder of this chapter.
First, the inductive biases of a GNN are encoded in the choice of aggregator and the edge-weighting scheme,
not in the overall template.
Second, the number of layers controls the receptive field
--- a node's representation after $L$ layers aggregates information from its $L$-hop neighbourhood ---
which, for skeleton graphs with at most four hops between any two joints,
places a natural upper bound on depth.

\subsection{Graph Convolutional Networks (GCN)}
\label{sec_meth_gcn}

The GCN architecture of~\cite{kipf2017gcn} derives from a first-order approximation of spectral graph convolutions.
Given adjacency matrix $A \in \mathbb{R}^{N \times N}$ and degree matrix $D = \mathrm{diag}(\sum_j A_{ij})$,
the symmetric normalised adjacency with self-loops is
\begin{equation}
    \tilde{A} = \tilde{D}^{-1/2}\,(A + I)\,\tilde{D}^{-1/2},
    \label{eq:gcn_adjacency}
\end{equation}
where $\tilde{D} = \mathrm{diag}(\sum_j (A+I)_{ij})$ is the degree matrix of the self-loop-augmented graph.
The GCN propagation rule is
\begin{equation}
    H^{(l+1)} = \sigma\!\left(\tilde{A}\,H^{(l)}\,W^{(l)}\right),
    \label{eq:gcn_layer}
\end{equation}
where $H^{(l)} \in \mathbb{R}^{N \times d_l}$ stacks the node features row-wise,
$W^{(l)} \in \mathbb{R}^{d_l \times d_{l+1}}$ is a learnable weight matrix,
and $\sigma(\cdot)$ is a pointwise nonlinearity.

Two properties of Equation~\eqref{eq:gcn_layer} motivate the move to attention-based aggregation
in the pose-grouping setting.
First, the edge weights are fixed:
once $\tilde{A}$ is constructed from the graph topology,
no learning signal modulates the contribution of individual neighbours.
Two neighbours of the same target node contribute in proportion
determined only by the degrees of the three nodes involved.
Second, this uniform treatment is informationally inadequate for skeleton graphs where edge semantics vary sharply:
the edge from the left shoulder to the left elbow encodes different information
from the edge from the left shoulder to the head,
and a pose-grouping embedding should be free to weight them differently.
Attention-based aggregation, introduced next, provides exactly this flexibility.

\subsection{Graph Attention Networks (GAT and GATv2)}
\label{sec_meth_gat}

GAT~\cite{velickovic2018gat} replaces the fixed normalisation of Equation~\eqref{eq:gcn_adjacency}
with a learned attention weight for each edge.
For a pair $(i,j)$ with $j \in \mathcal{N}(i)$, the unnormalised attention score is
\begin{equation}
    e_{ij}^{\mathrm{GAT}} = \mathrm{LeakyReLU}\!\left( \mathbf{a}^{\top} \left[\, W\mathbf{h}_i \,\|\, W\mathbf{h}_j \,\right] \right),
    \label{eq:gat_attn}
\end{equation}
where $W \in \mathbb{R}^{d' \times d}$ is a shared linear projection,
$\mathbf{a} \in \mathbb{R}^{2d'}$ is a learned attention vector,
and $[\cdot \| \cdot]$ denotes vector concatenation.
The scores are normalised over the neighbourhood with a softmax,
\begin{equation}
    \alpha_{ij} = \frac{\exp(e_{ij})}{\sum_{k \in \mathcal{N}(i)} \exp(e_{ik})},
    \label{eq:gat_softmax}
\end{equation}
and the updated node representation is the attention-weighted sum of projected neighbour features,
\begin{equation}
    \mathbf{h}_i' = \sigma\!\left( \sum_{j \in \mathcal{N}(i)} \alpha_{ij}\, W \mathbf{h}_j \right).
    \label{eq:gat_aggregate}
\end{equation}
Multi-head attention is recovered by repeating Equations~\eqref{eq:gat_attn}--\eqref{eq:gat_aggregate}
with independent parameters for each head and concatenating or averaging the outputs.

Equation~\eqref{eq:gat_attn} exposes a subtle limitation that was identified by~\cite{brody2022gatv2}
and has direct consequences for pose grouping.
Because $\mathbf{a}^{\top}$ is applied to the concatenation $[W\mathbf{h}_i \| W\mathbf{h}_j]$,
the attention score decomposes as $\mathbf{a}_1^{\top} W\mathbf{h}_i + \mathbf{a}_2^{\top} W\mathbf{h}_j$
where $\mathbf{a} = [\mathbf{a}_1; \mathbf{a}_2]$.
The term $\mathbf{a}_1^{\top} W \mathbf{h}_i$ is independent of $j$,
so the relative ordering of neighbours by attention score is the same for every query node
--- a property that~\cite{brody2022gatv2} formalises as \emph{static attention}.
In other words,
if node $j_1$ is ranked above node $j_2$ as a neighbour of $i$,
then $j_1$ is also ranked above $j_2$ as a neighbour of every other node.
For a pose-grouping network this is undesirable:
the network should be able to judge that a given left shoulder is important to a right shoulder
but not to a left ankle of a different person.

GATv2 resolves the limitation by reordering the operations inside the attention score,
\begin{equation}
    e_{ij}^{\mathrm{GATv2}} = \mathbf{a}^{\top}\, \mathrm{LeakyReLU}\!\left( W \left[\, \mathbf{h}_i \,\|\, \mathbf{h}_j \,\right] \right).
    \label{eq:gatv2_attn}
\end{equation}
With the nonlinearity now placed inside the query--key interaction,
the attention score does not decompose into independent per-node terms,
and the ranking of neighbours is allowed to depend on the query node.
GATv2 is strictly more expressive than GAT at the same parameter count~\cite{brody2022gatv2},
and every architecture evaluated in this chapter uses GATv2 as its attention primitive.
The remaining GCN and GAT references in Chapters~\ref{chap_4} and~\ref{chap_5}
refer to baselines and background, not to the working network.
